\chapter{Conclusiones y posibles ampliaciones}

\section{Cumplimiento de objetivos}

El trabajo entrega la infraestructura prometida en la introducción: una
plataforma sobre la que distintos modelos de decisión generados por la
primera fase pueden ejecutarse en un entorno común y producir, en cada
sesión, los cinco artefactos enumerados como objetivo general.

\begin{itemize}
\item La especificación del entorno se genera a partir de una
descripción en lenguaje natural y se valida estructuralmente antes de
ejecutarse.
\item El registro del comportamiento de los modelos se mantiene a lo
largo de la simulación con suficiente detalle para que sea analizable
posteriormente.
\item El análisis de patrones contrasta los datos observados con los
postulados disponibles en el Knowledge Backbone.
\item Los informes en PDF se generan automáticamente y contienen tanto
las métricas como las referencias a los paradigmas analizados.
\item Las observaciones de simulación se persisten como memorias
reutilizables, integrando la segunda fase como productora del backbone
compartido.
\end{itemize}

A esos artefactos hay que añadir un objetivo transversal igualmente
cumplido: la integración con la primera fase se sostiene sobre el
contrato \textit{duck-typed} que se describe a lo largo de la memoria,
y ha permitido que ambas fases evolucionaran de forma independiente sin
necesidad de coordinar publicaciones intermedias.

\section{Conclusiones técnicas}

Por encima de cada decisión concreta, la lección técnica que mejor
resume el trabajo es que la \emph{forma} en que se conectan los
componentes pesa más, en este sistema, que las elecciones internas de
cada uno. Las tres conclusiones que siguen son variaciones de esa
misma idea.

\paragraph{Contratos finos entre componentes que evolucionan a ritmos
distintos.} La integración con la primera fase podría haber requerido
un esquema cuidadosamente coordinado entre paquetes. En su lugar, se
redujo a un contrato \textit{duck-typed} mínimo, verificado en tiempo
de ejecución. Ese contrato es lo que ha hecho posible que las dos
fases del proyecto se desarrollaran en paralelo durante meses sin
coordinar publicaciones intermedias y sin pisarse en el esquema
compartido. La conclusión, en términos generales, es que la sencillez
del contrato pesa más que la rigidez de su declaración cuando las dos
partes están condenadas a evolucionar a ritmos distintos.

\paragraph{Centralizar el acceso al conocimiento merece la pena cuando
el coste de cada consulta no es despreciable.} La opción más obvia
habría sido dar a cada agente especializado acceso libre al
\textit{Knowledge Backbone}. Centralizar ese acceso en el Orchestrator,
con la única excepción del Analyst, dio un control fino sobre cuándo se
consulta el backbone y con qué presupuesto, y un único punto en el que
diagnosticar problemas de recuperación. La asimetría entre acceso
directo y acceso inyectado --- aparentemente complicada --- resultó más
útil en la práctica que un acceso uniforme: cada agente recibe el tipo
de contexto que de verdad necesita.

\paragraph{Separar la observación de la persistencia hace que el
backbone se mantenga limpio.} El Tracker se diseñó como observador
puro: produce un JSON estructurado y nada más. La escritura sobre el
grafo y los índices vectoriales ocurre después, a cargo del
\texttt{TrackerMemoryWriter} disparado por el Orchestrator. Un fallo en
el observador no deja memorias incompletas en el backbone, y el
\emph{writer} puede evolucionar (cambiar de embeddings, ajustar el
batch, cambiar el esquema) sin tocar el agente que observa. La lección
es la misma: definir bien la frontera entre las piezas vale más que
sofisticar cada una de ellas.

\section{Posibles ampliaciones}

El sistema deja varias líneas naturales de continuación.

\paragraph{Más paradigmas y otros tipos de entornos.} Cada nuevo modelo
que la primera fase produzca queda automáticamente disponible para
experimentar. A medio plazo, dar el salto a otros tipos de entorno más
allá de la rejilla 2D actual --- entornos continuos, escenarios
multi-agente con interacción cooperativa o competitiva, dominios no
espaciales --- permitiría estudiar paradigmas de toma de decisiones
para los que el escenario actual se queda corto.

\paragraph{Métricas cuantitativas estandarizadas.} La comparativa
multi-modelo actual combina métricas crudas con descripciones
cualitativas del Analyst. Una capa de métricas estandarizadas ---
recompensa acumulada, eficiencia, exploración, estabilidad de la
política --- haría las comparativas más directas y abriría la puerta a
un \textit{benchmark} reproducible para nuevos paradigmas.

\paragraph{Validación con el grupo de investigación y maduración del
producto.} El sistema se ha construido con el grupo de investigación
del tutor en mente, y el siguiente paso natural es ponerlo a su
disposición y recoger su valoración: qué paradigmas les interesa
estudiar, qué partes del informe les ahorran trabajo, qué echan en
falta. A partir de esa retroalimentación, la línea de maduración como
producto incluye una capa de gestión de usuarios y la posibilidad de
mantener varios hilos de conversación persistentes con el
Orchestrator, al estilo de un chat conversacional. A esto se
suman decisiones operativas que quedaron fuera del alcance de esta
fase --- copias de seguridad del Knowledge Backbone, monitorización ---
pero que cualquier despliegue real exigirá abordar.

\section{Una reflexión final}

Si la implementación deja de ser la actividad dominante, queda una
pregunta natural: ¿qué hacemos con ese margen recuperado? La respuesta
más interesante no es hacer lo mismo más deprisa, sino abordar sistemas
que antes habrían sido impracticables para un único autor en el plazo de
un trabajo de grado. Este propio TFG --- un pipeline de cinco agentes
sobre un backbone de cuatro componentes, dos fases coordinadas e
independientes y una interfaz web reactiva diseñada en Figma y
construida después con asistencia del agente --- difícilmente habría
alcanzado el estado en el que se entrega sin ese desplazamiento del
cuello de botella hacia las actividades de especificación, revisión y
verificación, como se describe en el capítulo de desarrollo.

Una vez aceptada esa nueva escala de lo posible, prescindir
voluntariamente de estas herramientas equivale a renunciar a parte del
alcance que permiten. El cambio de \textit{tooling} desplaza, en la
práctica, el techo de lo que un trabajo individual puede plantearse como
objetivo.

A modo de cierre, una observación más sobre el papel de la IA en este
trabajo. Tras meses construyendo el laboratorio en colaboración con un
agente, la imagen que queda es bastante distinta del relato de la IA
como amenaza al que es fácil asomarse. Lo que aporta una herramienta
así no es sustituir al desarrollador --- todas las decisiones de
alcance, las elecciones arquitectónicas y los criterios de aceptación
han seguido viniendo del lado humano ---, sino subir el nivel de
abstracción al que se trabaja. La pregunta deja de ser \emph{cómo lo
escribo} y pasa a ser \emph{qué quiero exactamente que ocurra}, lo que
permite plantear sistemas más ricos y más complejos sin renunciar a
controlarlos. Si esa lectura es correcta, lo razonable no es
defendernos de la IA, sino aprovecharla para hacer cosas que sin ella
no podríamos plantearnos --- exactamente lo que este TFG ha intentado
hacer.
